\documentclass{homework}
\usepackage[spanish]{babel}

%Paquetes AMS
\usepackage{amsmath}
\usepackage{amsfonts}
\usepackage{amssymb}

%Letra
%\usepackage{sans}
\usepackage{setspace}
%\spacing{1.15}

%Fondo
\usepackage[firstpage=true]{background}
% configuración
\backgroundsetup{
	scale=1.1, %escala de la imagen, es recomendable que sea del mismo tamaño que el pdf
	color=black, %fondo a usar para transparencia
	opacity=1.0, %nivel de transparencia
	angle=0, %en caso de querer una rotación
	contents={%
		\includegraphics[width=\paperwidth,height=\paperheight]{fondo.pdf} %nombre de la imagen a utilizar como fondo
	}%
}

%Márgenes, espacio y color
\usepackage{vmargin}
\usepackage{parskip}
\usepackage{setspace}
\usepackage{tcolorbox}
\usepackage{xcolor}

\setmargins{2.5cm}       % margen izquierdo
{2.0cm}                        % margen superior
{16.5cm}                      % anchura del texto
{24.0cm}                    % altura del texto
{5pt}                           % altura de los encabezados
{0.5cm}                           % espacio entre el texto y los encabezados
{0pt}                             % altura del pie de página
{1.5cm}                           % espacio entre el texto y el pie de página

%Encabezado y pie de página
\usepackage{fancyhdr}

\lhead{
	\begin{picture}(0,0) \put(0,0){\includegraphics[height=10mm]{logo-fciencias.jpeg}} \end{picture}
}

\rhead{Respuestas}
\renewcommand{\footrulewidth}{0.25pt}

\pagestyle{fancy}

% Gráficas y demás
\usepackage{pgfplots}
\usepackage{graphicx}

%\renewcommand{\baselinestretch}{1.08}

% Marcadores PDF
\usepackage{hyperref}
\hypersetup{
	colorlinks=true,
	linkcolor=blue,
	filecolor=magenta,      
	urlcolor=cyan,
	bookmarks=true,
}

\homeworksetup{
     username={Autores, por orden alfabético: \\ Shao Jie Hu Chen \\ Adrián Jaén Fuentes \\ Aarón Jerónimo Fernández \\ Noura Lachhab Bouhmadi \\ Laura Lázaro Soraluce \\ \vspace{1em} Doble Grado en Ingeniería Informática y Matemáticas},
     course={\huge{Estadística Descriptiva e Introducción a la Probabilidad \\ \vspace{1.5em} Relación de Ejercicios 3\vspace{1.5em}}}
}

%CUERPO
\begin{document}
	
	\spacing{1.5}
	\listofproblem
	\spacing{1.1}
	\newpage
	
%%%%%%%%%%%%%% EJERCICIO 1 %%%%%%%%%%%%%%
\problem
Se han lanzado dos dados varias veces, obteniendo los resultados que se presentan en la siguiente tabla, donde $X$ designa el resultado del primer dado  e $Y$ el resultado del segundo:
$$
   \begin{array}{|c|c|c|c|c|c|c|c|c|c|c|c|c|c|c|c|c|c|c|c|c|c|c|c|c|}
 X & 1 & 2 & 2 & 3 & 5 & 4 & 1 & 3 & 3 & 4 & 1 & 2 & 5 & 4 & 3 & 4 & 4 & 5 & 3 & 1 & 6 & 5 & 4 & 6 \\ \hline
 Y & 2 & 3 & 1 & 4 & 3 & 2 & 6 & 4 & 1 & 6 & 6 & 5 & 1 & 2 & 5 & 1 & 1 & 2 & 6 & 6 & 2 & 1 & 2 & 5 \\
   \end{array}
$$

  \begin{enumerate}
     \item Construir la tabla de frecuencias.
     \item Calcular las puntuaciones medias obtenidas con cada dado y ver cuales son m\'as homog\'eneas.
     \item ?`Qu\'e
      resultado del segundo dado es  m\'as frecuente cuando en el primero se obtiene
        un 3?
     \item Calcular la puntuaci\'on m\'axima del 50\% de las puntuaciones m\'as bajas obtenidas con el primer dado si con el segundo se ha obtenido un 2 o un 5.
   \end{enumerate}

Antes de resolver el ejercicio crearemos un nuevo enunciado, pues no tiene sentido calcular las cosas que nos piden, cuando estas dependen de la aleatoriedad. En nuestro nuevo enunciado, la variable X será el número de personas que componen distintas familias de Granada, y la variable Y es el número de habitaciones que tienen en sus casas. Resolvemos ahora el ejercicio.

\subproblem
Calculamos la tabla de frecuencia y además, los elementos que necesitaremos durante el ejercicio.

\begin{table}[h]
    \centering
    \begin{tabular}{|c|c|c|c|c|}
        \hline
         $x_i$ & $n_{i.}$ & $x_{i}n_{i.}$ & $x_i^2 n_{i.}$ \\ \hline
         1 & 4 & 4 & 4 \\ \hline 
         2 & 3 & 6 & 12 \\ \hline 
         3 & 5 & 15 & 45\\ \hline 
         4 & 6 & 24 & 96 \\ \hline 
         5 & 4 & 20 & 100 \\ \hline 
         6 & 2 & 12 & 72 \\ \hline 
          & 24 & 81 & 329 \\ \hline 
    \end{tabular}
\end{table}

\begin{table}[h]
    \centering
    \begin{tabular}{|c|c|c|c|}
        \hline
         $y_j$ & $n_{.j}$ & $y_{j}n_{.j}$ & $y_{j}^2n_{.j}$ \\ \hline
         1 & 6 & 6 & 6 \\ \hline 
         2 & 6 & 12 & 24 \\ \hline 
         3 & 2 & 6 & 18 \\ \hline 
         4 & 2 & 8 & 32 \\ \hline 
         5 & 3 & 15 & 75 \\ \hline 
         6 & 5 & 30 & 180 \\ \hline
           & 24 & 77 & 335 \\ \hline 
    \end{tabular}
\end{table}

\subproblem
Calculamos la media de nuestras variables:

\begin{equation*}
    \overline{x} = \dfrac{1}{n} \sum_{i=1}^k x_i n_{i.} = \dfrac{81}{24} = 3.375 \approx 3 hijos
    \hspace{1cm}
    \overline{y} = \dfrac{1}{n} \sum_{k=1}^p y_j n_{j.} = \dfrac{77}{24} = 3.2083 \approx 3 habitaciones
\end{equation*}

Ahora calculamos la desviación típica de cada variable para poder calcular el coeficiente de variación de Pearson y ver cuál es más homogénea:

\begin{equation*}
    \sigma_{x}^2 = \sqrt{\dfrac{328}{24}-3.375^2} = 1.5224 hijos
    \hspace{1cm}
    \sigma_{y} = \sqrt{\dfrac{335}{24}-3.2083^2} = 1.9144 habitaciones
\end{equation*}

\begin{equation*}
    CV_{x} = \dfrac{\sigma_x}{\overline{x}} = 0.4511
    \hspace{1cm}
    CV_{y} = \dfrac{\sigma_y}{\overline{y}} = 0.5967
\end{equation*}

Como vemos, los resultados de la primera variable son más homogéneos pues su coeficiente de variación de Pearson es menor.

\subproblem
Mirando la tabla de resultados, vemos que cuando las familias tenían 3 habitaciones en sus casas, el número de hijos más frecuente es el 4, con 2 veces, luego $M_0=4$.

\subproblem
Vamos a calcular la media de $X/Y$ cuando las familias tienen dos o cinco habitaciones en sus casa. Mirando la tabla de resultados obtenemos:

\begin{table}[ht]
    \begin{tabular}{|c|c|c|}
        \hline
         $x_i/Y$ & $n_{i}$ & $N_{i}$ \\ \hline
         1 & 1 & 1 \\ \hline 
         2 & 1 & 2 \\ \hline 
         3 & 1 & 3 \\ \hline 
         4 & 3 & 6 \\ \hline 
         5 & 1 & 7 \\ \hline  
         6 & 2 & 9 \\ \hline 
    \end{tabular}
\end{table}

Como $n/2 = 4.5$ y la frecuencia absoluta acumulada inmediatamente superior a este número es $N_{4} = 6$, deducimos que la mediana es $M_e = 4$.




%%%%%%%%%%%%%% EJERCICIO 2 %%%%%%%%%%%%%%
%PREAMBULO
\documentclass{article}
\usepackage{pgfplots}

%CUERPO
\begin{document}
    Ejercicio 2. La puntuación obtenida por 50 personas que se presentaron a una prueba de selección, sumadas las puntuaciones de los distintos tests fueron: 
    \\
    \\
    174, 185, 166, 176, 145, 166, 191, 175, 158, 156, 156, 187, 162, 172, 197, 181, 151,
    161, 183, 172, 162, 147, 178, 176, 141, 170, 171, 158, 184, 173, 169, 162, 172, 181,
    187, 177, 164, 171, 193, 183, 173, 179, 188, 179, 188, 179, 167, 178, 180, 168, 148, 173.
    \\
    \\
    POBLACIÓN: Las personas que se presentaron a la prueba de selección\\
    TAMAÑO: 50\\
    MODALIDADES: Los intervalos que contienen las notas que han obtenido las personas en la prueba
    
    \vspace{5mm}
    a) Agrupar los datos en intervalos de amplitud 5 desde 140 a 200 y dar la tabla de frecuencias.
    \\
    \begin{center}
    \begin{tabular}{| c | c | c | c | c |}
        \hline
        $x_i$ & $n_i$ & $N_i$ & $f_i$ & $F_i$ \\ \hline
        (140-145] & 2 & 2 & 0,04 & 0,04 \\
        (145-150] & 2 & 4 & 0,04 & 0,08 \\
        (150-155] & 1 & 5 & 0,02 & 0,1 \\
        (155-160] & 4 & 9 & 0,08 & 0,18 \\
        (160-165] & 5 & 14 & 0,1 & 0,28 \\
        (165-170] & 6 & 20 & 0,12 & 0,4 \\
        (170-175] & 10 & 30 & 0,2 & 0,6 \\
        (175-180] & 8 & 38 & 0,16 & 0,76 \\
        (180-185] & 6 & 44 & 0,12 & 0,88 \\
        (185-190] & 3 & 47 & 0,06 & 0,94 \\
        (190-195] & 2 & 49 & 0,04 & 0,98 \\
        (195-200] & 1 & 50 & 0,02 & 1 \\
        \hline
    \end{tabular} 
    \end{center}
    
    \vspace{5mm}
    
    Establecemos los intervalos de amplitud 5, empezando con 140 y terminando con 200. Para cada intervalo, contamos cuantas personas han obtenido una puntuación que esté dentro de dicho intervalo. Estas son las frecuencias absolutas ($n_i$). 
    Para la frecuencia absoluta acumulada ($N_i$), sumamos la frecuencia absoluta de la modalidad que estamos analizando a la frecuencia absoluta acumulada de la modalidad anterior. O lo que es lo mismo, sumamos las frecuencias absolutas de todas las modalidades hasta la que estamos analizando: $N_i=n_1 + n_2 + ... + n_i$. 
    Para frecuencia relativa, dividimos cada $n_i$ entre n, es decir, entre $N_{12}$, la frecuencia absoluta acumulada de la última modalidad. 
    Para la frecuencia relativa acumulada hacemos lo mismo que para la absoluta acumulada pero con la frecuencia relativa: sumamos las frecuencias relativas de todas las modalidades hasta la que estamos analizando: $F_i=f_1 + f_2 + ... + f_i$. 
    
    \vspace{10mm}
    b) Representar la distribución mediante un histograma, poligonal de frecuencias y curva de distribución.
    
    \begin{center}
    \begin{tikzpicture}
        \begin{axis}[
            title={Histograma},
            ymin=0, ymax=11,
            minor y tick num = 4,
            area style,
            ]
        \addplot+[ybar interval,mark=no] plot coordinates { (140, 2) (145, 2) (150, 1) (155, 4) (160, 5) (165, 6) (170, 10) (175, 8) (180, 6) (185, 3) (190, 2) (195, 1) (200, 0) };
        \end{axis}
    \end{tikzpicture}
    \end{center}
    
    Para el histograma, en el eje x tenemos las modalidades ($x_i$) y en el eje y un múltiplo de hi ($hi=\frac{fi}{a}$). En este caso, las amplitudes de todos los intervalos son iguales: 5. 
    
    \vspace{5mm}
    \begin{center}
    \begin{tikzpicture}
    \begin{axis}[
        title={Poligonal de frecuencias},
        xmin=135, xmax=205,
        ymin=0, ymax=11,
        xtick={140,150,160,170,180,190,200},
        ytick={0,2,4,6,8,10},
        legend pos=north west,
        ymajorgrids=false,
        grid style=dashed,
    ]

    \addplot[
        color=blue,
        ]
        coordinates {
        (142.5,2)(147.5,2)(152.5,1)(157.5,4)(162.5,5)(167.5,6)(172.5,10)(177.5,8)(182.5,6)(187.5,3)(192.5,2)(197.5,1)
        };
            [
            ymin=0, ymax=11,
            minor y tick num = 4,
            area style,
            ]
        \addplot+[ybar interval,mark=no] plot coordinates { (140, 2) (145, 2) (150, 1) (155, 4) (160, 5) (165, 6) (170, 10) (175, 8) (180, 6) (185, 3) (190, 2) (195, 1) (200, 0) };
    
    \end{axis}
    \end{tikzpicture}
    \end{center}
    
    Para el poligonal de frecuencias, los ejes son los mismos que para el histograma. En este caso, unimos los puntos que corresponden a las marcas de los intervalos en el histograma. Para hallar las marcas de los intervalos: $c_i=\frac{e_{i-1}+e_i}{2}$.
    
    \begin{center}
    \begin{tikzpicture}
    \begin{axis}[
        title={Curva de distribución},
        xmin=135, xmax=210,
        ymin=0, ymax=1.2,
        xtick={150,160,170,180,190,200},
        ytick={0,0.2,0.4,0.6,0.8,1},
        legend pos=north west,
        ymajorgrids=false,
        grid style=dashed,
    ]

    \addplot[
        color=blue,
        ]
        coordinates {
        (135,0)(140,0)(145,0.04)(150,0.08)(155,0.1)(160,0.18)(165,0.28)(170,0.4)(175,0.6)(180,0.76)(185,0.88)(190,0.94)(195,0.98)(200,1)(205,1)(210,1)
        };
    
    \end{axis}
    \end{tikzpicture}
    \end{center}
    
    En la curva de distribución, en el eje x tenemos los extremos superiores de los intervalos $(e_i)$ y en el eje y, $F(e_i)$. Pintamos la curva que une los puntos tal que $F(e_i)=\displaystyle\sum_{j=1}^{i}f_{j}=F_i$. En este caso tenemos que la curva es continua. $F(e_i)$ es 0 para valores menores que $e_1$ y 1 para valores mayores que el último intervalo $e_k$.
    
\end{document}




%%%%%%%%%%%%%% EJERCICIO 3 %%%%%%%%%%%%%%
\subproblem
Completar la tabla
\begin{center}
%	\begin{table}[htbp]
		\begin{center}
			\begin{tabular}{|c|c|c|c|c|c|c|c|}
				\hline
				$I$ & $n_{i}$ & $N_{i}$ & $f_{i}$ & $F_{i}$ & $c_{i}$ & $a{i}$ & $h_{i}$\\\hline
				(8300, 9300] & 2 & 2 & 2/18 & 2/18 & 8800 & 1000 & 0.002/18 \\ \hline
				(9300, 10200] & 3 & 5 & 3/18 & 5/18 & 9750 & 900 & 1/5400 \\ \hline
				(10200, 11300] & 5 & 10 & 5/18 & 10/18 & 10750 & 1100 & 1/3960 \\ \hline
				(11300, 12700] & 2 & 12 & 2/18 & 12/18 & 12000 & 1400 & 1/12600 \\ \hline
				(12700, 13500] & 4 & 16 & 4/18 & 16/18 & 13100 & 800 & 0.005/18 \\ \hline
				(13500, 14500] & 2 & 18 & 2/18 & 1 & 14000 & 1000 & 0.002/18\\ \hline
			\end{tabular}
		\end{center}
%	\end{table}
\end{center}
\subproblem
Representar la distribución mediante un histograma, poligonal de frecuencias y curva de
distribución.

\begin{center}
En el histograma se representan los $I_{i}$ en el eje X y reescalando $h_{i}$ multipliplicándolo por 10.000, para poder tener una mejor representación en el eje Y. 
\end{center}

\begin{center}
    \begin{tikzpicture}
        \begin{axis}[
            title = {Histograma},
            xmin=7000, xmax=15000,
            area style
            ]
        \addplot+[ybar interval,mark=no] plot coordinates {
        (8300, 1.111)
        (9300, 1.852)
        (10200, 2.525)
        (11300, 0.794)
        (12700, 2.777)
        (13500, 1.111)
        (14500, 0)
        };
        \end{axis}
    \end{tikzpicture}
\end{center}

\begin{center}
En la poligonal de frecuencias unimos los "techos", es decir las marcas de clase $a_{i}$, de cada intervalo $I_{i}$.
\end{center}

\vspace{5mm}
\begin{center}
    \begin{tikzpicture}
        \begin{axis}[
        title={Poligonal de frecuencias},
        xmin=7000, xmax=15000,
        ymin=0
        ]

        \addplot[
        color=blue,
        ]
        coordinates {
        (8800, 1.111)
        (9750, 1.852)
        (10750, 2.525)
        (12000, 0.794)
        (13100, 2.777)
        (14000, 1.111)
        };
        [
        ymin=0, ymax=11,
        minor y tick num = 4,
        area style,
        ]
        \addplot+[ybar interval,mark=no] plot coordinates {
        (8300, 1.111)
        (9300, 1.852)
        (10200, 2.525)
        (11300, 0.794)
        (12700, 2.777)
        (13500, 1.111)
        (14500, 1.111)
        };
        \end{axis}
    \end{tikzpicture}
\end{center}

\begin{center}
En la curva de distribución representamos los $I_{i}$ en el eje X y las frecuencias absolutas acumuladas $N_{i}$ en el eje Y.
\end{center}

\begin{center}
    \begin{tikzpicture}
        \begin{axis}[
        title={Curva de distribución},
        xmin=7000, xmax=15500,
        ymin=-0.1, ymax=1.2,
        %xtick={1/18, 5/18, 10/18, 12/18, 16/18, 1},
        %ytick={0,0.2,0.4,0.6,0.8,1},
        legend pos=north west,
        ymajorgrids=false,
        grid style=dashed,
        ]

        \addplot[
        color=blue,
        ]
        coordinates {
        (700,0)
        (8300,0)
        (9300, 2/18)
        (10200, 5/18)
        (11300,10/18)
        (12700,12/18)
        (13500,16/18)
        (14500,1)
        (15500,1)
        }; 
        \end{axis}
    \end{tikzpicture}
\end{center} 

\subproblem
¿Cuántas comunidades presentan una renta menor o igual que 12700 euros? ¿Y cuántas superior a 11300 euros?

Como podemos ver en la tabla, las comunidades que presentan una renta menor o igual a 12700 corresponden con $N_{4}$, es decir, 12 comunidades.
Las comunidades con una renta superior a 11300 corresponderían con la suma de las comunidades de los intervalos superiores, es decir, a la suma de $n_{4}$, $n_{5}$ y $n_{6}$. Finalmente, el resultado a esta última pregunta es 8 comunidades.




%%%%%%%%%%%%%% EJERCICIO 4 %%%%%%%%%%%%%%
\problem
Sea $X$ una variable aleatoria con funci{\'o}n de densidad
$$
f(x) = \left \{
\begin{array}{lc}
k_{1} (x+1) & 0 \leq x \leq 4 \\
\\ k_{2} x^{2} & 4 <x \leq 6
\end{array}
\right.
$$
Sabiendo que  $P(0 \leq X \leq 4) = 2/3$,  determinar $k_{1},\ k_{2}$, y
deducir su funci{\'o}n de distribuci{\'o}n.

Nótese que la variable aleatoria es continua y toma valores entre 0 y 6. En el desarrollo de este ejercicio, se presupondrá que la función de masa dada es correcta. 

En primer lugar, calculemos la función de distribución de la variable aleatoria. Aplicamos la siguiente expresión: $F(x) = \int_{-\infty}^x f(t) dt$. 

Si $0 \leq x \leq 4$, tenemos que: 

$$F(x) = \int_{-\infty}^x f(t) dt = \int_{-\infty}^0 0 dt + \int_0^x k_1(t+1)dt = k_1(\frac{·x^2}{2} + x)$$ 

Si $4 < x \leq 6$, tenemos que:

\begin{equation*}
\begin{split}
F(x) & = \int_{-\infty}^x f(t) dt = \int_{-\infty}^0 0 dt + \int_0^4 k_1(t+1)dt + \int_4^x k_2·t^2 dt \\
& = k_1(\frac{4^2}{2} + 4) + k_2(\frac{x^3}{3} - \frac{4^3}{3}) 
\end{split}
\end{equation*}

Sabemos que $P(0\leq X \leq 4) = F(4) - F(0^-)$. Como la variable es continua, tenemos que $F(0^-) = F(0)$. Deducimos entonces que: 

$$P(0 \leq X \leq 4) = F(4) - F(0^-) = F(4) - F(0) = k_1(\frac{4^2}{2} + 4) = 2/3 \Rightarrow k_1 = \frac{1}{18}$$

Por otra parte, para que se verifique que se trate, en efecto, de una función de distribución, se ha de verificar que: 

$$F(+\infty) = k_2(\frac{6^3}{3} - \frac{4^3}{3}) + \frac{2}{3} = 1 \Rightarrow k_2 = \frac{1}{152}$$

Por tanto, la función de distribución viene dado por la siguiente expresión: 

\begin{equation*}
	F(x) = \left \{
	\begin{array}{lcc}
	0 & x < 0 \\
	\frac{1}{18}(\frac{x^2}{2} + x) & 0 \leq x \leq 4 \\
	\frac{x^3}{456} + \frac{10}{19} & 4 <x \leq 6 \\
	1 & x > 6
	\end{array}
	\right.
\end{equation*}




%%%%%%%%%%%%%% EJERCICIO 5 %%%%%%%%%%%%%%
Escribamos las tablas con los datos que nos pueden interesar para la resolución del ejercicio: \\

\begin{tabular}{| c | c | c | c | c |}
	\hline
	$I_i^{(1)}$ & $c_i^{(1)}$ & $n_i^{(1)}$ & $N_i^{(1)}$ & $h_{i} = \frac{n_{i}}{a_{i}}$ \\ \hline
	$(0, 1]$ & $0.5$ & $12$ & $12$ & $12$ \\
	$(1, 2]$ & $1.5$ & $13$ & $25$ & $13$  \\
	$(2, 3]$ & $2.5$ & $11$ & $36$ & $11$  \\
	$(3, 4]$ & $3.5$ & $8$ & $44$  & $8$   \\
	$(4, 5]$ & $4.5$ & $6$ & $50$  & $6$ \\ \hline
	
\end{tabular} \\\\\\

\begin{tabular}{| c | c | c | c | c |}
	\hline
	$I_i^{(2)}$ & $c_i^{(2)}$ & $n_i^{(2)}$ & $N_i^{(2)} $ & $h_{i} = \frac{n_{i}}{a_{i}}$ \\ \hline
	$(0, 1]$ & $0.5$ & $1$ & $1$ & $1$ \\
	$(1, 3]$ & $2$ & $6$ & $7$ & $3$ \\
	$(3, 6]$ & $4.5$ & $7$ & $14$ & $2.333$ \\
	$(6, 10]$ & $8$ & $12$ & $26$ & $3$ \\
	$(10, 12]$ & $11$ & $2$ & $28$ & $1$ \\ \hline
	
\end{tabular}\\
\\ 
\\
\subproblem
Medias aritm{\'e}tica, arm{\'o}nica y geom{\'e}trica.

La media aritmética de una variable es la suma de sus valores entre en número total de observaciones. Como los datos están organizados en intervalos de clase, para calcular la media aritmética vamos a suponer que todos los datos de un intervalo son idénticos a la marca de clase de cada intervalo. Por tanto, la media la calculamos de la siguiente manera: 
	
	$$\overline{x} = \dfrac{1}{n}\sum_{i=1}^{k}n_i·c_i$$


Para la primera distribución, $n = 50$, luego la media aritmética es $2.16$ u. 

Para la segunda, $n = 28$, luego la media aritmética será 5.786 u. \\


La media armónica se usa para promediar datos de magnitudes relativas. La definimos como la inversa de la media aritmética de los valores inversos de la variable (en nuestro caso, usamos las marcas de clase):
\begin{center}
	$H = \dfrac{n}{\dfrac{n_1}{c_1}+\dfrac{n_2}{c_2}+ ...+\dfrac{n_k}{c_k}} = \dfrac{n}{\sum_{i=1}^{k}\dfrac{n_i}{c_i}}$
\end{center}

Para la primera distribución: H = 1.229 u.

Para la segunda distribución: H = 3.399 u.  \\


Finalmente la media geométrica se usa cuando se desea promediar datos de una variable que tiene efectos
multiplicativos acumulativos en la evolución de una determinada característica con un
valor inicial fijo.

Es la raíz n-ésima del producto de los $n$ valores (o marcas de clase) de la distribución:

\begin{center}
	$G = \sqrt[n]{\prod_{i=1}^{k}c_i^{n_i}}$
\end{center}

Para no perder precisión en el resultado, la calcularemos sabiendo que el logaritmo de la media geométrica es la media aritmética de los logaritmos de los valores de la variable:

\begin{center}
	$\log G = \log \sqrt[n]{\prod_{i=1}^{k}c_i^{n_i}} = \dfrac{1}{n}\sum_{i=1}^{k}n_i \log c_i$
\end{center}

Para la primera distribución: G = 1.685 u.

Para la segunda distribución: G = 4,769 u.  \\


\subproblem

Se nos pide calcular la moda de las distribuciones, es decir, el valor  de mayor frecuencia: \\

En el caso primero, observamos que se encuentra en el intervalo (1,2], luego haremos un promedio teniendo en cuenta los intervalos contiguos. Para ello apliquemos la interpolación lineal (pues estamos suponiendo que los datos están distribuidos uniformemente), además de suponer que es una distribución continua.

\begin{center}
	$M_O = e_{i-1} + \dfrac{h_i - h_{i-1}}{2h_i - h_{i-1} - h_{i+1}} (e_i - e_{i-1})$
\end{center}

Por tanto, la moda será: $M_O = 1.\overline{3}$.\\


En el segundo caso, tenemos dos modas, pues el valor máximo de $h_{i}$ lo toman dos intervalos. De esta forma, deducimos que existen dos modas: la primera en el intervalo (1,3] y otra en el intervalo (6,10]. En el primer caso, la moda se calcula con expresión anterior empleando $e_{i-1} = 1, e_i=3, h_{i-1} = 1, h_i = 3, h_{i+1} = 2.333 $, deducimos que $M_O = 2.5 u$. En el segundo, que proviene del intervalo (6, 10], aplicamos el mismo método teniendo en cuenta que $e_{i-1} = 6, e_i=10, h_{i-1} = 2.333, h_i = 3, h_{i+1} = 2 $, $M_O = 7 u$.

\subproblem

Para calcular la mediana hemos de observar las frecuencias absolutas acumuladas; tendremos que encontrar el punto que divida a la población en dos partes iguales, es decir, $n/2$:

Para la primera distribución: $n/2 = 25$ luego el $50\%$ de la población supera el valor 2.

Para la segunda distribución: $n/2 = 14$ luego el $50\%$ de la población supera el valor 6.\\
\\

\subproblem

\begin{enumerate}
	\item Recorrido: Es la amplitud del intervalo en la que se mueven los valores de la variable, por tanto, se calcula restando el valor mas grande posible menos el más pequeño. Como en nuestro caso tenemos intervalos, cojamos el extremo superior del último intervalo de clase y le restamos el extremo inferior del primer intervalo de clase:
	
	Para la primera distribución será 5 y para la segunda será 12.
	\item Recorrido intercuartílico: Es la magnitud que indica la longitud del intervalo que contiene al $50\%$ central de los datos observados. Para calcularlo necesitamos los cuartiles 1 y 3:
	
	\begin{center}
		$Qk = e_{i-1} + \dfrac{\dfrac{n}{k} - N_{i-1}}{n_i}(e_i-e_{i-1})$
	\end{center}
	
	Para la distribución 1 será $Q_1 = 1.038 u, Q_3 = 3.187 u; R_I = 2.1 u$
	
	Para la distribución 2 será $Q_1 = 3 u, Q_3 = 8.333 u; R_I = 5.333 u$\\
	
	\item Desviación típica: es una medida de dispersión, una medida de cómo los valores individuales de el conjunto pueden diferir de la media. Es la raíz cuadrada de la media de las distancias al cuadrado de cada uno de los valores de la variable a la media de la distribución. También se puede calcular como la raíz cuadrada positiva de la varianza.
	\begin{center}
		$\sigma^2 = \dfrac{1}{n} \sum_{i=1}^{n}(x_i-\overline{x})^2$
	\end{center}
	
	Para la distribución 1 será $\sigma^2 = 1.744; \sigma = 1.321$
	
	Para la distribución 1 será $\sigma^2 = 8.522; \sigma = 2.919$\\
	
	
	Para ver cuál es mas homogénea, calculemos el Coeficiente de variación de Pearson:
	
	$C.V._1 = \dfrac{1.321}{2.16}= 0.612$ \\
	\\
	\\
	$C.V._2 = \dfrac{2.919}{5.786}= 0.504$
	
	La segunda distribución es más homogénea.
	
\end{enumerate}



%%%%%%%%%%%%%% EJERCICIO 6 %%%%%%%%%%%%%%
\problem

Para cada una de las distribuciones:

\hskip 2cm Distribuci{\'o}n A  \hskip 2cm Distribuci{\'o}n B \hskip 2.2cm Distribuci{\'o}n C
\vskip -0.5cm $$
\left .
\begin{array}{c|ccc}
X \backslash Y & 10 & 15 & 20 \\ \hline
1  &  0 &  2 &  0  \\
2  &  1 &  0 &  0  \\
3  &  0 &  0 &  3  \\
4  &  0 &  1 &  0  \\
\end{array}
\right .
\hskip 1cm
\left .
\begin{array}{c|ccc}
X \backslash Y & 10 & 15 & 20 \\ \hline
1  &  0 &  2 &  0  \\
2  &  1 &  0 &  0  \\
3  &  0 &  0 &  3  \\
\end{array}
\right .
\hskip 1cm
\left .
\begin{array}{c|cccc}
X \backslash Y & 10 & 15 & 20  & 25 \\ \hline
1  &  0 &  3 &  0  & 1 \\
2  &  0 &  0 &  1  & 0 \\
3  &  2 &  0 &  0  & 0 \\
\end{array}
\right .
$$

\begin{enumerate}
	\item ?`Dependen funcionalmente $X$ de $Y$ o $Y$ de $X$?
	\item Calcular las curvas de regresi{\'o}n y comentar los resultados.
	
\end{enumerate}

\begin{center}
    \begin{tabular}{|c|c c c c|c|c|c|c|}
        \hline
        $X/Y$ & 1 & 2 & 3 & 4 & $n_i.$ & $n_i.\cdot x_i $ & $n_i.$ \cdot $x_i^2$ & $x_i\sum_{j=1}^p n_{ij} y_j$\\ \hline
        10    & 1 & 3 & 0 & 0 & 4 & 40 & 400 &  70\\
        12    & 0 & 1 & 4 & 3 & 8 & 96 & 1152 & 312\\
        14    & 2 & 0 & 0 & 2 & 4 & 56 & 784 &  140\\
        16    & 4 & 0 & 0 & 0 & 4 & 64 & 1024 & 64\\ \hline
        $n._j$ & 7 & 4 & 4 & 5 & 20 & 256 & 3360 & 586 \\ \hline
        $n._j\cdot y_j$ & 7 & 8 & 12 & 20 & 47 & & &   \\ \hline
        $n._j\cdot y_j^2$ & 7 & 16 & 36 & 80 & 139 & & &  \\ \hline
    \end{tabular}
\end{center}
   
\\
\subproblem
¿Son estadísticamente independientes $X$ e $Y$?\\

Para que se diera la independencia estadística, se tendría que cumplir que las frecuencias de $X$ condicionadas a $Y$ son las mismas para valor de $x_i$. 
$$ \dfrac{n_{i1}}{n._1}=\dfrac{n_{i2}}{n._2}= \dots = \dfrac{n_{ij}}{n._j} = \dots = \dfrac{n_{ip}}{n._p} \qquad \forall i = 1,2,\dots , k.$$
Pero como en algunos casos $n_{ij} = 0$, no se puede dar la independencia ya que no todos son nulos.\\

\subproblem Calcular y representar las curvas de regresión de $X/Y$ e $Y/X$.\\

La curva de regresión tipo 1 de X sobre Y vendrá dada por los puntos $(x_i, \overline{y_i}), i = 1, \dots k$.

$$\{(10, 1.75), (12, 3.25), (14, 2.5), (16, 1)\}$$\\


La curva de regresión tipo 1 de Y sobre X vendrá dada por los puntos $(y_j, \overline{x_j}), i = 1, \dots k$.

$$\{(1, 14.571), (2, 10.5), (3, 12), (4, 12.8) \}$$\\

\subproblem
Cuantificar el grado en que cada variable es explicada por la otra mediante la correspondiente curva de regresión. \\

Con los datos de la tabla:

$$\overline{x} = \dfrac{256}{20} = 12.8 \qquad \overline{y} = \dfrac{27}{20} = 2.35 $$

$$ \sigma_x^2 = m_{20} - m_{10}^2= 4.16 \qquad \sigma_y^2 = m_{02} - m_{01}^2= 1.4275 $$ \\

Calculemos la varianza de los residuos para $Y/X$:
$$\sigma_{ry}^2 = \sum_{i=1}^{k}\sum_{j=1}^p n_{ij}\cdot(y_j-f(x_i))^2 = \dfrac{13.25}{20} = 0.6625 $$

Para $X/Y$:
$$\sigma_{ry}^2 = \sum_{i=1}^{k}\sum_{j=1}^p n_{ij}\cdot(x_i-f(y_j))^2 = \dfrac{37.51}{20} = 1.875 $$

Podemos ahora calcular el coeficiente de correlación:

$$\eta_{\frac{y}{x}}^2 = 1 - \dfrac{\sigma_{ry}}{\sigma_y^2} = 0.536 \qquad 
\eta_{\frac{x}{y}}^2 = 1 - \dfrac{\sigma_{rx}}{\sigma_x^2} = 0.549 $$

Con lo que podemos ver que la bondad está por debajo del 55\%, por lo que la curva no es demasiado precisa. \\
\\

\subproblem 
¿Están X e Y correladas linealmente? Dar las expresiones de las rectas de regresión.

Calculemos el coeficiente de correlación lineal:

$$r =  \dfrac{\sigma_{xy}}{\sigma_x \sigma_y} = -0.32$$

Podemos observar una baja correlación lineal.\\

Calculamos las rectas:

Y sobre X : 
$$y = ax + b$$

$$a = \dfrac{\sigma_{xy}}{\sigma_x^2} = -0.187 \qquad b = \overline{y} - a\overline{x} = 4.75 $$

Recta de Y sobre X $\equiv y = -0.187x + 4.75$\\

X sobre Y :

$$x = ay + b$$

$$a = \dfrac{\sigma_{xy}}{\sigma_y^2} = -0.55 \qquad b = \overline{x} - a\overline{y} = 14 $$

Recta de X sobre Y $\equiv x = -0.55y + 14$ 




%%%%%%%%%%%%%% EJERCICIO 7 %%%%%%%%%%%%%%
\subproblem
A = Sacar blanca en la primera urna

B = Sacar blanca en la segunda urna

C = Sacar blanca en la tercera urna \\


$P(Sacar\ 4 \ blancas) = P(U_1)P(A|U_1) + P(U_2)P(B|U_2) + P(U_3)P(C|U_3) = \\


\dfrac{1}{3}(\dfrac{5}{10}\cdot\dfrac{4}{9}\cdot\dfrac{3}{8}\cdot\dfrac{2}{7}) + 
  \dfrac{1}{3}(\dfrac{6}{10}\cdot\dfrac{5}{9}\cdot\dfrac{4}{8}\cdot\dfrac{3}{7}) + 
  \dfrac{1}{3}(\dfrac{7}{10}\cdot\dfrac{6}{9}\cdot\dfrac{5}{8}\cdot\dfrac{4}{7}) = 0.0873$ \\

\subproblem

$A_i$ = Coger una urna i 

$B$ = Sacar 4 y que solo una sea negra \\

$P(A_i) = \dfrac{1}{3}$

$P(B|A_1) = \dfrac{5}{10} \cdot \dfrac{4}{9} \cdot \dfrac{3}{8} \cdot \dfrac{5}{7} = \dfrac{5}{84} = 0.0595$ \\

$P(B|A_2) = \dfrac{6}{10} \cdot \dfrac{5}{9} \cdot \dfrac{4}{8} \cdot \dfrac{4}{7} = \dfrac{2}{21} = 0.0952$ \\

$P(B|A_3) = \dfrac{7}{10} \cdot \dfrac{6}{9} \cdot \dfrac{5}{8} \cdot \dfrac{3}{7} = \dfrac{1}{8} = 0.125$ \\


Para calcular $P(A_2|B)$ usamos la regla de Bayes:

$$ P(A_2|B) = \dfrac{P(B|A_2)P(A_2)}{P(B|A_1)P(A_1)+P(B|A_2)P(A_2)+P(B|A_3)P(A_2)} = \dfrac{\frac{2}{21 \cdot 3}}{\frac{1}{3}(\frac{5}{84} + \frac{2}{21} + \frac{1}{8})} = \dfrac{16}{47}$$




%%%%%%%%%%%%%% EJERCICIO 8 %%%%%%%%%%%%%%
\problem
	 De una muestra de 24 puestos de  venta en  un  mercado  de
	abastos  se ha recogido informaci\'on sobre el n\'umero de balanzas ($X$) y el n\'umero de dependientes ($Y$). Los resultados aparecen en la siguiente tabla:
	$$
	\begin{array}{c|cccc}
		X \backslash Y & 1  &  2 &  3 & 4\\ \hline
		1  &  1 &  2 &  0 & 0 \\
		2  &  1 &  2 &  3 & 1 \\
		3  &  0 &  1 &  2 & 6 \\
		4  &  0 &  0 &  2 & 3 \\
	\end{array}
	$$
	\begin{enumerate}
		\item Determinar las rectas de regresi{\'o}n.
		\item ?`Es apropiado suponer que existe una relaci{\'o}n  lineal  entre  las
		variables?
		\item  Predecir, a partir de los resultados, el n\'umero de balanzas que puede esperarse en un puesto con seis dependientes.  ?`Es fiable esta predicci{\'o}n?
	\end{enumerate}

	En una población de tamaño n = 24 se ha observado dos variables estadísticas,
	X = número de balanzas e Y = número de dependientes, las cuales han
	presentado k = 4, p = 4 modalidades distintas, con distribución de frecuencia
	conjunta $(x_{i},y_{j})$, $n_{ij}$ $ i = 1,...,3$ $ j=1,...3$
	
	$\bullet$ Empezamos rellenando nuestra tabla: \\
	
	\begin{tabular}{ | c | c | c | c | c | c | c | c| c| }
		
		
		\hline	
		X \ Y & 1 & 2 & 3 & 4  & $n_{i.}$ & $n_{i.}x_{i.}$ &  $n_{i.}x_{i.}^{2}$ & $x_{i.}\sum_{j=i}^{p}n_{ij}y_{j}$ \\ \hline
		1 & 1 & 2  & 0 & 0 & 3 & 3 & 3 & 5 \\
		2 & 1 & 2 & 3 & 1& 7 & 14 & 28 & 36 \\
		3 & 0 & 1 & 2 & 6 & 9 & 27 & 81 & 96 \\
		4& 0 & 0  & 2  & 3 & 5 & 20 & 80  & 72 \\
		$n_{.j}$& 2 & 5  & 7 & 10 & 24 & 64 & 192 & 209 \\ 
		$n_{.j}y_{.j}$ & 2 & 10 & 21 & 40 & 73 & & &  \\
		$n_{.j}y_{.j}^{2}$& 2 & 20 & 63 & 160 & 245 & & & \\\hline
		
		
	\end{tabular}
	\\ \\
	
\subproblem	
	$\bullet$ Para ello primero calculamos las medias de X e Y: \\
	
	$\bar{x} = \frac{3 + 14 + 27 +20}{24} = \frac{64}{24} = 2,6667 $ balanzas
	\\ 
	
	$\bar{y} = \frac{2 + 10+21 +40}{24} = \frac{73}{24} = 3,0417  $ dependientes
	\\
	\\
	
	$\bullet$ Tenemos que calcular también las varianzas y la covarianza: \\
	
	
	
	$\sigma_{x}^{2} =  $$\frac{1}{n}\sum^{k}_{i=1}n_{i.}x_{i}^{2} - \bar{x}^{2} = \frac{192}{24} - 2,6667^{2} = 0,8887 $ $ balanzas^{2} $ \\
	
	$\sigma_{y}^{2} = $ $\frac{1}{n}\sum^{p}_{j=1}n_{.j}y_{j}^{2} - \bar{y}^{2} = \frac{245}{24} - 3,0417 ^{2} = 0,9564 $ $dependientes^{2}$ \\
	
	$\sigma_{xy} = $ $\frac{1}{n}\sum^{k}_{i=1}\sum^{p}_{j=1}n_{ij}x_{i}y_{j} - \bar{xy}= \frac{1}{24}.209 - 8,113 = 0,5970 $ \\
	
	
	$\bullet$Ahora que tenemos todos los datos necesarios podemos ya calcular los coeficientes de las rectas de regresión :\\ \\
	\textbf{Para Y / X:}\\
	$y = ax +b $\\ \\
	$a = \frac{\sigma_{xy}}{\sigma_{x}^{2}} = \frac{0,5970}{0,8887} = 0,6718$\\ \\
	$b = \bar{y} - a.\bar{x} = 3,0417 - \frac{0,5970}{0,8887}. 2,6667 = 1,2502$\\ \\
	La recta de regresión de Y sobre X es
	$y = 0,6718 + 1,2502 $\\ \\
	\textbf{Para X / Y:}\\
	$x = ay + b $\\ \\
	$a = \frac{\sigma_{xy}}{\sigma_{y}^{2}} = \frac{0,5970}{ 0,9564 } = 0,6242$\\ \\
	$b = \bar{x} - a.\bar{y} = 2,6667 -  \frac{0,5970}{ 0,9564 }.3,0417 = 0,768$\\ \\La recta de regresión de X sobre Y es
	$x = 0,6242y +  0,768 $\\ 
\subproblem	
	
	
	$\bullet$ Para ello tenemos que calcular el coeficiente de correlación lineal:\\ \\
	
	$r^{2} = \frac{\sigma_{xy}^{2}}{\sigma_{x}^{2}\sigma_{y}^{2}} = \frac{0,5970^{2}}{0,8887 . 0,9564} = 0,4193$\\ \\
 
		$\bullet$El resultado que hemos obtenido nos indica que la recta de regresión de Y sobre X nos da una información de menos
	del 42 \% de la variabilidad de Y, luego no seria buena idea suponer una relación lineal entre las variables ya que la bondad de la función es relativamente baja. 
	Tampoco, el coeficiente de correlación lineal $r =\sqrt{r^{2}} = 0,6475$ nos indica que existe una buena correlación lineal directa entre las variables. \\ \\
	
\subproblem	
	$x = 0,6242y +  0,786 = 0,6242 . 6 +  0,786 = 4,5312 \approx 5 $ balanzas. Si bien la recta de regresión proporcionaría una estimación de la cantidad de balanzas que cabría esperar para 6 dependiente, el resultado no sería fiable, puesto que la recta de regresión tendría únicamente validez para el intervalo muestral estudiado (en este caso, de 1 a 4 dependientes). Más allá de estos límites, no se sabe si la nube de puntos seguiría comportándose como los datos estudiados. 




%%%%%%%%%%%%%% EJERCICIO 9 %%%%%%%%%%%%%%
\problem
\ vskip 0.75cm \ item Fonctions calculées   de la densité de las variables $ Y = 2
X + 3 $  \ y \  $ Z = | X | $ , siendo   $ X $ una variable continua con funci { \ ' o} n de
densité
$$
f_X (x) =   \ frac {1} {4}, \ \ - 2 <x < 2 .
$$
\subproblem\\ \\	
\begin{flushleft}
	
		Sea X, una variable aleatoria continua, con f(x) su función de densidad y E=(-2,2).
\end{flushleft}
\\
	Como Y=h(X), entonces $h^{-1}(y) = X = \frac{Y-3}{2} $\\ \\
	Calculamos la imagen de h(x) en los extremos del intervalo E=(-2,2):\\
	$h(-2)= -1 $\\
	$h(2)= 7 $
	\\ \\
Por lo tanto $E^{'} = (-1, 7)$\\ \\
Sea:\\
 $g(y) = f(h^{-1}(y)) . |(h^{-1})^{'}(y)|$, si $y\in (-1,7)$\\
$g(y) = 0 $ en otro caso.\\ \\
Calculamos g(y):\\
$g(y) = f(\frac{Y-3}{2}) . |\frac{1}{2}| = \frac{1}{4}.\frac{1}{2} = \frac{1}{8}$ ,  si $y\in (-1,7)$ \\ \\
Veamos ahora la variable Z. Como Z=h(X), entonces $h^{-1}(z) = X $\\ \\
Z=$|X|$=\left\lbrace\begin{array}{c} -x~si~x~<=~0\\  x~si~x~>~0 \end{array}\right.\\

\subproblem\\ \\
\begin{flushleft}
	Calculamos la imagen de h(E) en los extremos del intervalo E=(0,2):
\end{flushleft}
\\ \\
$g(y) = f(h^{-1}(y)) . |(h^{-1})^{'}(y)| + f(h^{-1}(y)) . |(h^{-1})^{'}(y)| = \frac{1}{4}.1 + \frac{1}{4}.1 = \frac{1}{2} $, si\\ $z\in(0,2)$\\
\\ \\
 g(y)=\left\lbrace\begin{array}{c} \frac{1}{2}~~si~0~<~x~<~2\\ 0~~~en~otro~caso \end{array}\right.\\






%%%%%%%%%%%%%% EJERCICIO 10 %%%%%%%%%%%%%%
Escribamos las tablas con los datos que nos pueden interesar para la resolución del ejercicio: \\

\begin{tabular}{| c | c | c | c | c | c |}
	\hline
	$I_i$   & $n_i$ &$N_i$&$c_i$&$h_i$ & $n_i(x_i-\overline{x})^2$ \\ \hline
	$(10, 30]$ & 15 & 15  & 20  & 0.75 & 12355.55\\
	$(30, 40]$ & 22 & 37  & 35  & 2.2  & 4129.18 \\
	$(40, 50]$ & 48 & 85  & 45  & 4.8  & 657.12  \\
	$(50, 60]$ & 40 & 125 & 55  & 4    & 1587.6  \\
	$(60, 90]$ & 25 & 150 & 75  & 0.833& 17292.25\\ \hline
	           & 150&     &     &      & 36021.7 \\ \hline
	
\end{tabular} \\\\\\

\subproblem

Se nos pide calcular la edad más común, luego buscamos la moda de la distribución. Dado que los intervalos son de diferente tamaño tenemos que calcular las densidades de frecuencia de cada intervalo. El mayor $h_i$ es 4.8 luego la moda está en el intervalo (40, 50]. Para calcularla tendremos que dibujar el histograma con el intervalo en cuestión y los dos contiguos. Unimos la esquina superior derecha del primer rectángulo con la del segundo, y la esquina superior izquierda del segundo con la del tercero. Ahora, aplicando semejanza de triángulos calculamos la intersección, con lo que obtenemos la moda.
\begin{center}
	$\dfrac{4.8-2.2}{4.8-4} = \dfrac{M_O-40}{50-M_O}; M_O = 47.647 \approx 48$ años
\end{center}


\subproblem

Se nos pide calcular el percentil 35 y el 65. Hay que buscar el intervalo asociado al valor inmediatamente superior a $n*p$ en la frecuencia absoluta acumulada. Después, se calcula aplicando semejanza de triángulos : \\


\begin{center}
	$n*0.35 = 52.5; N_i = 85; I_i = (40, 50]$\\
	$\dfrac{52.5-37}{P_{35}-40} = \dfrac{85-37}{50-40} = 43.229
	\approx 43$ años.
\end{center}

Luego la edad mínima será aproximadamente 43 años.

\begin{center}
	$n*0.65 = 97.5; N_i = 125 ; I_i = (50, 60]$\\
	$\dfrac{97.5-85}{P_{65}-50} = \dfrac{125-85}{60-50} = 53.125
	\approx 53$ años.
\end{center}

Por tanto, la edad máxima será aproximadamente 53 años.\\


\subproblem

Para el recorrido intercuartílico necesitaremos los cuartiles 1 y 3, es decir, los percentiles 25 y 75. Procedemos de la misma manera que en el apartado anterior:\\ \\ 

\begin{center}
	$n*0.35 = 37.5\quad  N_i = 85\quad  I_i = (40, 50]$\\
\end{center}

\begin{center}
	$\dfrac{37.5-37}{Q_{1}-40} = \dfrac{85-37}{50-40} = 40.104
	\approx 40$ años.
\end{center}

\begin{center}
	$n*0.75 = 112.5\quad  N_i = 125 \quad  I_i = (50, 60]$\\
	
\end{center}

\begin{center}
	$\dfrac{112.5-85}{Q_{3}-50} = \dfrac{125-85}{60-50} = 56.875
	\approx 57$ años.
\end{center}

Por tanto el recorrido intercuartílico será $Q_3-Q_1 = 16.771 \approx 17$ años, lo que nos indica que el 50\% de la población central se encuentra en un intervalo de unos 17 años.\\

Para la calcular la desviación típica necesitamos calcular la raíz cuadrada positiva de la varianza. Tenemos que calcular la media del área de los cuadrados de lado la diferencia de cada dato a la media. La media es 48.7. Escribamos los $n_i(x_i-\overline{x})^2$ en la tabla. La varianza será 240.144 años$^2$. Finalmente la desviación típica será 15.497.


\subproblem

\begin{enumerate}
	\item Coeficiente de asimetría de Fisher:
	$$\gamma_1(X) = \dfrac{\mu_3}{\sigma_X^3};\qquad
	\mu_3 = m_3 - 3m_2m_1 + 2m_1^3 = \dfrac{1}{n}\sum_{i} n_i ci^3 - \dfrac{3}{n}
	\sum_{i} n_i ci^2\overline{x}  + 2\overline{x}^3 = 341.256$$
    
    $\gamma_1(X) = \dfrac{341.256}{15.4965^3} = 0.0917 > 0 \longrightarrow $ Asimetría moderada por la derecha.

	\item Coeficientes de asimetría de Pearson:
	Necesitamos la mediana para calcular A*. 150*0.5 = 75 $\in I_3$
	$$Me = 40 + \dfrac{75-37}{48}*10 = 47.91\overline{6}$$

	$$A_p = \dfrac{\overline{x}- M_O}{\sigma_X} = 0.06795 \qquad A_p^* = \dfrac{3(\overline{x}- Me)}{\sigma_X} = 0.1516$$
    
    Nos indican también una ligera asimetría por la derecha.

	\item Coeficiente de curtosis de Fisher:

	$$\gamma_2(X) = \dfrac{\mu_4}{\sigma^4} - 3$$
    
    $$\mu_4 = m_4 - 4m_3m_1 + 6m_2m_1^2 - 2m_1^4 = 153232.46$$
    
    $$\gamma_2(X) = \dfrac{153232.46}{15.4965^2} - 3 = -0.3429 < 0$$
	
	Por tanto, la distribución es platicúrtica; presenta menor concentración central de frecuencias que una distribución normal con su media y desviación típica,
	
	\item Coeficiente de curtosis de Kelley:
	
	$$K = \dfrac{1}{2} \dfrac{Q_3 - Q_1}{D_9 - D_1}- 0,263$$
	
	Sabemos $Q_1$ y $Q_2$, luego calcularemos $D_1$ y $D_9$
	
	$1n\alpha = 15 \longrightarrow D_1$ coincide con el extremo superior de $I_1$, luego será 30 años.
	$1n\alpha = 135 \longrightarrow D_9 \in I_5$. Aplicamos la fórmula.
	\begin{center}
	    $ D_9 = 60 + \dfrac{135- 125}{25}*30 = 72 $ años
	\end{center}
	 $$ K = -0.06335 < 0 $$
Como $K < 0$, la distribución es platicúrtica.
\end{enumerate}


	
\end{document}